\chapter{Historia del desarrollo (II): de la selección secuencial a la cartera}
\label{chap:desarrollo-cartera}

Esta segunda mitad retoma la historia del desarrollo iniciada en el
Capítulo~\ref{chap:desarrollo-metodo}. Si aquella cubría el paso a un único Model Study y los
defectos que afectaron a la metodología de contratos y a la selección secuencial, ésta cubre el
rediseño de la construcción de cartera descrita en el Capítulo~\ref{chap:diseno-experimental}: tres
defectos con una raíz común en la lógica de venta, el colapso de la calibración isotónica que
motiva la salvaguarda de la curva de alfa usada en la Sección~\ref{sec:cartera-barrido}, y dos
decisiones tomadas contra la recomendación técnica. Cierra con la advertencia de reproducibilidad
que condiciona la lectura de toda cifra del Capítulo~\ref{chap:resultados}.

\section{Toda venta necesita un destino mejor}

El rediseño de la cartera partió de tres defectos que, vistos por separado, parecían independientes,
y que compartían una única raíz: las ventas se decidían sin mirar a dónde iba el dinero.

El primero era el más grave en términos contables. Bajo la política de efectivo oportunista, si
ningún candidato superaba el umbral, la lógica vendía la cartera entera y devolvía una cartera
objetivo vacía. El \textit{backtest} interpretaba «vacía» como «mantener» y resucitaba las
posiciones --- pero ya había cargado los costes de unas ventas que nunca se ejecutaron. La cartera
pagaba por replegarse sin replegarse.

El segundo producía el efecto contrario y contribuyó de forma decisiva al 877\,\% de rotación anual
que se observaba entonces. Bajo la política de permanecer siempre invertido, las doce posiciones se
vendían por umbral y el relleno obligatorio \emph{recompraba exactamente las mismas en el mismo
snapshot}: una ida y vuelta completa, con sus costes, para acabar donde se estaba.

El tercero era la ausencia de un límite de concentración: con un único candidato admisible y un tope
de efectivo del 20\,\%, el 80\,\% de la cartera acababa en un solo valor.

La regla que los resuelve a los tres es la que gobierna hoy la cartera: una venta sólo se emite si
el destino del dinero es mejor que la posición, después de costes. El destino puede ser otro valor
--- y entonces debe superar el umbral de rotación --- o el efectivo, que exige un tope declarado y
respeta un suelo de diversificación. Se añadió además histéresis entre el umbral de entrada y el de
salida, para que una posición no oscile dentro y fuera por variaciones mínimas.

\section{El colapso de la calibración isotónica}

El diagnóstico de posiciones que nunca se vendían pese a un percentil muy bajo, presentado en el
Capítulo~\ref{chap:desarrollo-metodo}, encadenaba dos problemas distintos. Aquel capítulo documentó
el primero, de naturaleza predictiva --- una ventana rectangular que no olvidaba el régimen
antiguo. Éste es el segundo, más sutil, y bloqueaba la cartera por completo.

La calibración de percentil a alfa esperada usaba una regresión isotónica creciente. Cuando la
relación real entre percentil y retorno futuro es \emph{decreciente}, la única curva creciente que
minimiza el error es una constante. El resultado fue que la alfa esperada era idéntica para los
quinientos cuatro valores del universo. Con una alfa plana no existe ventaja calculable entre dos
posiciones, y el bucle de rotación se interrumpía sin emitir ninguna orden. No era un fallo de
implementación: era el comportamiento correcto de un estimador aplicado fuera de su supuesto.

El diagnóstico de ese episodio contiene una lección estadística que merece registrarse por sí misma.
El primer análisis por deciles, hecho con la \emph{media} del retorno excedente, sugería que la
ordenación estaba invertida en todo el histórico. Era un artefacto del estimador: la distribución
tiene una cola derecha extrema, y unos pocos valores multiplicadores situados en el decil bajo
bastaban para que su media (${+}3552$ puntos básicos) contradijera a su mediana (${-}146$). Con la
mediana, o con una media winsorizada, la curva histórica \emph{sí} es creciente, de forma coherente
con el Rank-IC positivo. La inversión era real, pero sólo en el régimen reciente. Un diagnóstico
apresurado habría llevado a reescribir el sistema entero para corregir un problema que no existía.

La isotónica se sustituyó por una recta de percentil a retorno anualizado real con una cascada de
ventanas, que acepta la primera ventana con pendiente creciente. La estimación se hace sobre veinte
ventiles, pero la evaluación se realiza en el rango continuo: produce quinientos cuatro valores de
alfa distintos donde la isotónica producía uno. Es la cascada que el
Capítulo~\ref{chap:diseno-experimental} describe en «De rango a alfa esperado».

\section{Dos decisiones tomadas contra la recomendación técnica}

El registro documenta dos decisiones que se adoptaron en contra de la recomendación técnica. Se
recogen aquí con ambos lados del argumento, porque omitirlas presentaría el protocolo como más
limpio de lo que es.

La primera fue sustituir la calibración isotónica sin conservarla como opción del catálogo. A favor:
el colapso es demostrable y deja la cartera inoperativa, de modo que mantener la opción habría
significado conservar una rama que se sabe rota. En contra: cambia el ganador sin pasar por la
selección secuencial por Rank-IC, que es precisamente la regla 4 del protocolo, y renuncia a poder
comparar ambas formulaciones con evidencia. El autor asumió el riesgo y pidió eliminar la
alternativa.

La segunda es más incómoda. Cuando ninguna de las tres ventanas de la cascada produce una pendiente
creciente, el sistema aplica una salvaguarda: una recta fija entre $-10\,\%$ y $+10\,\%$. En contra
de esa elección se argumentó que es un supuesto \emph{a priori}, del mismo tipo que el sesgo de
régimen fijado por constantes que este proyecto ya había eliminado, y que se activa precisamente
cuando las tres ventanas coinciden en que la ordenación no discrimina a favor --- es decir, cuando
la evidencia disponible dice que no hay señal, el sistema inventa una. La recomendación técnica fue
usar una curva plana, equivalente a repartir por igual. Se prefirió la recta fija.

La mitigación consiste en registrar, para cada fila, qué ventana produjo su alfa. En el estudio de
referencia la salvaguarda se activa en 4.384 de 20.545 filas, algo más de una de cada cinco. Es una
proporción suficientemente alta como para que el lector deba tenerla presente al interpretar
cualquier cifra de cartera, y por eso se recoge también en el Capítulo~\ref{chap:limitaciones}.

Durante esa misma revisión se descubrió que la ruta de calibración de alfa \emph{no tenía ninguna
cobertura de pruebas}. Un renombrado había dejado una llamada rota en esa ruta que la batería
existente no detectaba. Se añadió un contrato específico que cubre justamente lo que puede fallar en
silencio: que cada rango produzca una alfa distinta y monótona, que una relación decreciente caiga a
la salvaguarda, y que la ausencia de cohortes cerradas produzca un valor no disponible en lugar de
un cero.

\section{Una posición sin puntuación detenía toda la rotación}

Un episodio concreto ilustra por qué las casuísticas de datos incompletos merecen tratamiento
explícito. En una ejecución, una posición estaba en el percentil 4,96 con un retorno esperado
de $-9{,}01\,\%$ anual, mientras otro valor fuera de la cartera estaba en el percentil 100 con
$+10{,}00\,\%$: una ventaja de 1.901 puntos básicos, muy por encima de cualquier umbral de rotación.
El intercambio no se produjo.

La causa era que una tercera posición, de una empresa que había salido del índice años antes, no
tenía fila puntuable en esa fecha. El bucle la seleccionaba como peor posición, no podía calcular su
ventaja y se detenía --- bloqueando de paso todas las rotaciones posteriores, incluida la que sí era
claramente rentable. La corrección fue tratar la ausencia de puntuación como un mandato de venta,
no como una decisión económica: una posición sin puntuación actual se vende, ignorando el mínimo de
tenencia, y no se recompra en ese mismo snapshot.

\section{Advertencia de reproducibilidad}
\label{sec:reproducibilidad}

Hay una consecuencia de este historial que debe declararse sin rodeos. El estudio del que procede
toda la evidencia de este trabajo, \studyid, se ejecutó bajo la versión 5 del catálogo. El código
actual está en la versión 6, tras la corrección del reparto de efectivo del 4 de agosto de 2026.

Los artefactos siguen siendo válidos y trazables: conservan sus identificadores, sus hashes de
conjunto de datos y su instantánea de catálogo, y cualquier cifra de este documento puede
verificarse contra ellos. Pero \textbf{no son reproducibles bit a bit con el código actual}, y no
son comparables con estudios ejecutados bajo la versión 6.

La alternativa habría sido reejecutar el estudio completo. Se descartó por una razón metodológica y
no de conveniencia: la reejecución cambiaría el ganador, y escribir los capítulos de resultados
sobre un ganador y después sustituirlo por otro cuando el texto ya está redactado se parece
demasiado a elegir la configuración cuyos resultados encajan mejor con lo escrito. Se optó por
congelar el estudio, declarar la discrepancia y mantener todas las cifras referidas a la versión con
la que realmente se calcularon.
